\chapter{12}



Inmitten vereinzelter Bäume entdeckte er eine Bank, steuerte sie an und liess sich
erleichtert darauf niederfallen. Er hatte freie Sicht auf die Baustelle. Konzentriert betrachtete er die 
Umgebung. Still war es, förmlich ausgestorben. Am Feierabend der
Bauarbeiter konnte es wohl kaum liegen. Dafür war es zu früh. Vielleicht
war hier ein Baustopp verordnet worden? 
Da und dort öffnete sich in der Ferne die Tür eines Wohnhauses. Jemand 
ging hinein oder kam heraus.
Louis zog den Rucksack näher an sich heran und tastete darin nach
der Taschenlampe und dem Päckchen Batterien. Er versuchte, sich
einzig auf das Hochschieben des Kunststoffplättchens zu konzentrieren
und legte zwei Stück ein. Dann sass er da, verwünschte die drückende Hitze und dachte nach. Er würde die Dunkelheit abwarten, in das leere Gebäude einsteigen und am frühen Morgen wieder verschwinden. Bis dahin mussten erstmal einige Stunden vergehen. 
Mit einem Ächzen verstaute er die Einkaufstüte im Rucksack, als er im Augenwinkel eine Frau mit Hund dem Weg entlangkommen sah.
Mit einem "Grüezi" zogen die beiden an ihm vorbei. Der Gruss verblüffte ihn ebenso, wie er ihn beruhigte. 
Sie musste ihn als einen gewöhnlichen, jungen Mann wahrgenommen haben. Augenblicklich wusste er, was zu tun war. 
Er stand entschlossen auf, warf sich den Rucksack über die Schulter und zog los. Der Zürichsee hatte ihm schon einmal das Warten erleichtert. 

\sectionbreak

Der Einstieg über die Ostseite des Rohbaus ging leicht. Die noch
fensterlosen Räume wirkten dumpf. Wände mit Löchern. Es war
stockfinster.
Als er sich sicher fühlte, zog er leise seine Taschenlampe aus dem
Rucksack. Er leuchtete die dunkeln Ecken des Raumes aus und richtete
sich in einer geschützten Nische zum Schlafen ein. Es war staubig, der
Boden uneben und in der Luft hing der Geruch nach feuchtem Beton. Also
doch. Hier wurde gebaut.
Er atmete flach. Bewegte sich vorsichtig. In der Hoffnung, niemand
käme auf die Idee, die Baustelle nachts zu kontrollieren, konnte er sich
endlich fallen lassen. Jetzt war die Wärme angenehm. Frieren musste er
nicht, nicht auch noch. Er schlief augenblicklich ein.

\sectionbreak

Bis plötzlich - Louis schrak schweissgebadet hoch, unfähig erst, zu
unterscheiden zwischen Traum und Wirklichkeit. Er riss seine Augen auf
und starrte fassungslos in das grelle Licht einer Taschenlampe. Er
stiess einen Schrei in den Bau. Der Widerhall war gewaltig.
Sie hatten ihn.
